\section{Experimental protocols and data lineage}
\label{app:experiments}

All paths below are relative to the repository root.  The numerical kernels use
float64 arrays.  The objective reported throughout is the Gaussian variational
objective, up to an additive target-dependent constant,
\[
  \calE(m,C)
  =-\frac12\log\det C+\E_{X\sim\N(m,C)}V(X).
\]
When deterministic quadrature is used, NumPy's Gauss--Hermite nodes
\((x_j,w_j)\) are rescaled according to
\[
  \E f(Z)=\frac1{\sqrt\pi}\sum_{j=1}^{K}w_j f(\sqrt2x_j),
  \qquad Z\sim\N(0,1).
\]
Unless stated otherwise, a run is evolved for its fixed iteration budget; a
``time to tolerance'' is a diagnostic computed from the saved trajectory, not
an early-stopping rule.  The PDF files used in \Cref{sec:numerics} are
byte-for-byte copies of the corresponding files in
\path{reports/assets/figs}.  The script
\path{reports/make_report_assets.py} performs only post-processing and never
reruns a dynamics experiment.

\subsection{Covariance bootstrap}

Let
\[
  H=\diag(\nu_1,\ldots,\nu_d),
  \qquad (\nu_i)_{i=1}^d=\operatorname{logspace}(1,\kappa,d).
\]
The Gaussian target is
\[
  V(x)=\frac12x^\top Hx,
  \qquad (m_\star,C_\star)=(0,H^{-1}),
\]
and the non-Gaussian target used only in the transport experiment is
\[
  V(x)=\frac12x^\top Hx+\gamma\sum_{i=1}^d\log\cosh(x_i),
  \qquad \gamma=1.
\]
For the latter, \(m_\star=0\), and each diagonal entry of \(C_\star\) solves
\[
  c_i^{-1}=\nu_i+\gamma\E_{Z\sim\N(0,c_i)}\operatorname{sech}^2 Z.
\]
The fixed-point tolerance is \(10^{-13}\), with at most 500 iterations and
80-point Gauss--Hermite quadrature.  Thus \(\alpha=1\), while
\(\beta=\kappa\) for the Gaussian target and \(\beta=\kappa+1\) for the
\(\log\cosh\) target.

Write \(G=\E\nabla V(X)\), \(A=\E\nabla^2V(X)\).  The two Fisher--Rao steps
share
\[
  m_{n+1}=m_n-hC_nG_n
\]
and use, respectively,
\begin{align*}
  C_{n+1}^{\Riem}
  &=C_n^{1/2}\exp\!\left[h\{\Id-C_n^{1/2}A_nC_n^{1/2}\}\right]C_n^{1/2},\\
  C_{n+1}^{\KL}
  &=(1+h)(C_n^{-1}+hA_n)^{-1}.
\end{align*}
Every run starts from
\[
  m_0=m_\star+3\sqrt{\diag(C_\star)},
  \qquad C_0=\lambda_0\Id.
\]

The Gaussian scaling grid is
\[
  d\in\{10,20\},\quad
  \kappa\in\{10,100,1000\},\quad
  \lambda_0\in\{10^{-10},10^{-8},10^{-6},10^{-4},10^{-2}\},
\]
for both covariance maps, with \(h=0.5\) and 200 steps.  The burn-in time is
the first \(n\) such that \(\lambda_{\min}(C_n)\geq(2\beta)^{-1}\).
The left panel of \Cref{fig:num-three-stage} selects KL, \(d=20\),
\(\kappa=1000\), and \(\lambda_0=10^{-10}\).  Its three ordinates are
\(\lambda_{\min}(C_n)\), the exact Gaussian energy gap, and
\(\norm{C_n^{1/2}G_n}\).  Its vertical markers are the covariance threshold
and the first saved iterate with energy gap at most \(10^{-3}\).

For the transport grid, the Gaussian and \(\log\cosh\) targets both have
\(d=20\) and \(\kappa=50\).  It uses
\[
  \lambda_0\in\{10^{-8},10^{-6},10^{-4}\},\qquad
  c\in\{0.25,0.5,1\},\qquad h=0.5,
\]
both Fisher--Rao tails, and a 6000-step budget.  When \(\lambda_0<c/\beta\),
one Bures step of size \(\eta=c/\beta\) is applied:
\begin{align*}
  m^+&=m-\eta G,\qquad
  \widetilde C=(\Id-\eta A)C(\Id-\eta A),\\
  C^+&=\frac12\left\{
    \widetilde C+2\eta\Id+
    [\widetilde C(\widetilde C+4\eta\Id)]^{1/2}
  \right\}.
\end{align*}
The right panel of \Cref{fig:num-three-stage} selects \(c=0.5\) and plots,
for every target, tail, and \(\lambda_0\), the first iteration whose gap is at
most \(10^{-1}\), together with \(\lambda_{\min}(C)\) after the first step.

\paragraph{Data and figure lineage.}
\begingroup\small\raggedright
The configuration is
\path{configs/natural_gradient_covariance_bootstrap/covariance_bootstrap.yaml};
the runner is
\path{scripts/natural_gradient_covariance_bootstrap/run_experiments.py}; and
the kernels and figure selectors are in
\path{src/natural_gradient_covariance_bootstrap}.  The left panel reads
\path{outputs/natural_gradient_covariance_bootstrap/results_long.csv} through
\path{plotting.py::fig_three_stage}.  The right panel reads
\path{outputs/natural_gradient_covariance_bootstrap/wasserstein_bootstrap_summary.csv}
through \path{plotting.py::fig_wasserstein_bootstrap}.  The final assets are
\path{reports/assets/figs/fig_covboot_three_stage.pdf} and
\path{reports/assets/figs/fig_covboot_wasserstein_bootstrap.pdf}.
\par\endgroup

\subsection{Deterministic discretization grid}

All three main targets have dimension two and use
\(\lambda\in\{0.01,0.1,1\}\).  They are
\begin{align*}
  V_{\rm G}(x,y)
  &=\frac12(x,y)\diag(1,\lambda)(x,y)^\top,\\
  V_{\rm q}(x,y)
  &=\frac{(\sqrt\lambda x-y)^2}{20}+\frac{y^4}{20},\\
  V_{\rm s}(x,y)
  &=\frac{(\sqrt\lambda x-y)^2}{20}
    +\frac{\delta}{2}(x^2+y^2)+\gamma\log\cosh y,
  \qquad (\delta,\gamma)=(0.05,1).
\end{align*}
The Gaussian initialization is
\(m_0=(10,10)\), \(C_0=\diag(0.5,2)\); the other two use
\(m_0=(10,10)\), \(C_0=4\Id\).  Gaussian and quartic expectations are analytic,
using exact Gaussian moments through order four.  The smooth target uses
80-point one-dimensional Gauss--Hermite quadrature for the \(y\)-marginal.
The Gaussian optimum is analytic.  The other optima are computed in Cholesky
coordinates by four deterministic L-BFGS-B starts, with a 5000-iteration
budget, gradient tolerance \(10^{-10}\), and function tolerance \(10^{-16}\).

The common grid is
\[
  h\in\{0.001,0.002,0.005,0.01,0.02,0.05,0.1,0.2,
           0.5,1,2,5,10\},
  \qquad T=15,\qquad N=\lceil T/h\rceil.
\]
The steps are the two formulas in the preceding subsection, with exact
population \(G_n,A_n\).  A Radau solution of the continuous flow, with relative
tolerance \(10^{-10}\), absolute tolerance \(10^{-12}\), and 400 output
times, is used only for the terminal discretization-accuracy diagnostic.
Simulation stops early only on a nonfinite update or
\(\lambda_{\min}(C_n)\leq10^{-12}\).  A completed trajectory is called stable
when it is SPD-feasible, ends below its initial energy, and never exceeds
\(10^3\) times its initial energy gap.  Monotonicity allows an absolute
per-step tolerance of \(10^{-10}\); accuracy means
\[
  \norm{m_N-m(T)}+
  \frac{\norm{C_N-C(T)}_{\Frob}}{\norm{C(T)}_{\Frob}}
  \leq10^{-2}.
\]

The left panel of \Cref{fig:num-discretization} uses only \(V_{\rm s}\), all
three \(\lambda\)'s, and the matched stable steps
\[
  h\in\{0.05,0.1,0.5,1\}.
\]
It plots energy gap against \(nh\).  The right
panel retains every stable run that reaches gap \(10^{-6}\).  For each target,
method, and \(h\), it averages the hitting iteration and hitting time \(nh\)
over the eligible \(\lambda\)'s.  Its top row is iteration count and its bottom
row is physical flow time, not measured wall time.

\paragraph{Data and figure lineage.}
\begingroup\small\raggedright
The configuration and runner are
\path{configs/natural_gradient_discretization_stepsize/stepsize_grid.yaml} and
\path{scripts/natural_gradient_discretization_stepsize/run_stepsize_grid.py}.
Target, method, optimizer, ODE, and classification code is in
\path{src/natural_gradient_discretization_stepsize}.  Both selected panels read
\path{outputs/natural_gradient_discretization_stepsize/results_long.csv} and/or
\path{outputs/natural_gradient_discretization_stepsize/summary.csv}; their
builders are
\path{scripts/natural_gradient_discretization_stepsize/plot_results.py::fig_convergence_speed_grid}
and \path{plot_results.py::fig_iterations_vs_stepsize}.  The assets are
\path{reports/assets/figs/fig_convergence_speed_smooth_logconcave.pdf} and
\path{reports/assets/figs/fig_iterations_vs_stepsize.pdf}.
\par\endgroup

\subsection{STL estimator variance and constant-step floors}

Here \(A=\diag(a_1,\ldots,a_d)\), with \((a_i)\) log-spaced in
\([1,\kappa]\).  The targets are
\begin{align*}
  V_{\rm G}(x)&=\frac12x^\top Ax,\\
  V_{\rm lc}(x)&=\frac12x^\top Ax+\tau\sum_{i=1}^d\log\cosh(x_i).
\end{align*}
The grid is
\begin{align*}
  V_{\rm G}:&\quad d\in\{2,8,32\},\quad
       \kappa\in\{10,100,10^4\},\\
  V_{\rm lc}:&\quad d\in\{2,8,32\},\quad
       \kappa\in\{10,100\},\quad \tau\in\{0.1,1\}.
\end{align*}
For \(V_{\rm lc}\), \(m_\star=0\) and
\[
  (C_\star)_{ii}^{-1}
  =a_i+\tau\E_{Z\sim\N(0,(C_\star)_{ii})}\operatorname{sech}^2Z.
\]
This fixed point uses 80-point Gauss--Hermite quadrature, tolerance
\(10^{-13}\), and at most 500 iterations.

For the estimator-level experiment,
\[
  b_{\rm base}(X)=\nabla\log\pi(X),\qquad
  b_{\STL}(X)=\nabla\log\pi(X)+C^{-1}(X-m).
\]
At each target configuration the four states plotted on the left of
\Cref{fig:num-stl} have \(C=C_\star\) and
\[
  m=m_\star+\rho\sqrt{\diag(C_\star)},
  \qquad \rho\in\{0,0.25,1,4\}.
\]
The full estimator grid additionally includes
\((m_\star,0.25C_\star)\) and \((m_\star,4C_\star)\).  For every state and
each seed \(s\in\{0,\ldots,19\}\), \(M=100000\) Gaussian samples are drawn
in one chunk (the configured chunk cap is 200000).  The reported mean-block
Fisher--Rao variance is
\[
  \Var_{\FR}(b)=
  \E\bigl[(b-\E b)^\top C(b-\E b)\bigr].
\]
The left panel plots the median over the 20 seeds of
\(\Var_{\FR}(b_{\STL})/\Var_{\FR}(b_{\rm base})\) against \(\rho\), separately
for all 12 non-Gaussian configurations.  For display on a logarithmic axis,
the plotting code adds \(10^{-3}\) to \(\rho\) and clips ratios below
\(10^{-16}\).

The algorithm-level experiment compares
\(\{\Riem,\Riem\text{-STL},\KL,\KL\text{-STL}\}\).  At every step it draws
\(B\) samples from the current Gaussian, averages the corresponding mean
estimator, and averages the diagonal pointwise log-target Hessian for the
covariance step.  The two mean estimators within a covariance scheme use the
same standard-normal stream.  The settings are
\[
  h=0.1,\qquad B\in\{1,4,16\},\qquad
  m_0=m_\star+3\sqrt{\diag(C_\star)},\qquad C_0=0.5C_\star,
\]
with 2000 steps for \(V_{\rm G}\) and 3000 for \(V_{\rm lc}\).  Gaussian
energy gaps are analytic; non-Gaussian gaps use 40-point Gauss--Hermite
quadrature and are therefore independent of the mini-batch used for the
update.  An emergency eigenvalue guard clips values below \(10^{-12}\).  The
runner counts such clips in a cell diagnostic printed during execution, but the
surviving summary CSV does not preserve that count.

The replicate labels are \(0,\ldots,19\).  For the algorithm grid, one
batched random stream advances these 20 leading-array entries.  Its cell seed
is the first 15 hexadecimal digits of
\[
  \operatorname{SHA256}(
  \texttt{20260623|kind|d|kappa|tau|B}),
\]
interpreted as an integer; the method name is deliberately absent, so all four
paired methods receive the same standard-normal sequence.  For the estimator
grid, the generator seed is the replicate label itself.

For each run the empirical floor is the median gap over the final 20 percent
of the full, undecimated trajectory.  The right panel of \Cref{fig:num-stl}
then takes the median of these per-run medians over every target configuration
and seed, separately by target kind, method, and \(B\).  Nonpositive display
values are clipped to \(10^{-16}\).

\paragraph{Data and figure lineage.}
\begingroup\small\raggedright
The configuration is
\path{configs/natural_gradient_stl_variance/stl_variance.yaml}.  The two
runners are
\path{scripts/natural_gradient_stl_variance/run_estimator_variance.py} and
\path{scripts/natural_gradient_stl_variance/run_algorithm_grid.py}; kernels and
plot selectors are in \path{src/natural_gradient_stl_variance}.  The left panel
reads \path{outputs/natural_gradient_stl_variance/estimator_variance.csv} through
\path{plotting.py::fig_variance_by_distance} with
\texttt{kind=log\_cosh}.  The right panel reads
\path{outputs/natural_gradient_stl_variance/tail_noise_floor.csv} through
\path{plotting.py::fig_batchsize_effect}.  The assets are
\path{reports/assets/figs/fig_stl_variance_by_distance_to_optimum.pdf} and
\path{reports/assets/figs/fig_stl_batchsize_effect.pdf}.
\par\endgroup

\subsection{Trace and traceless affine modes}

For the Gaussian calibration target, \(\pi=\N(0,\Id_d)\), expectations are
exact.  The covariance step is diagonalized as
\[
  \lambda_i^+
  =\lambda_i\exp\!\left[
     \frac{h}{2\omega}\{-\lambda_i+q_\tau(C)\}
  \right],
  \qquad
  q_\tau(C)=\frac{\omega+\tau\Tr C}{\omega+d\tau},
\]
and \(m^+=m-hCm\).  The grid is
\[
  d\in\{2,5,10\},\quad
  \omega\in\{0.125,0.25,0.5,1,2\},\quad
  \tau\in\{-\omega/(2d),0,\omega/(2d)\},
\]
with \(h=0.02\), \(T=20\), and metrics saved every five steps.

The non-Gaussian target is
\[
  V_\rho(x)=\frac12\norm{x}^2+
  \frac\rho M\sum_{\ell=1}^{M}\log\cosh(a_\ell^\top x),
  \qquad d=5,\quad \rho=5,\quad M=4d=20.
\]
The unit vectors \(a_\ell\) are normalized Gaussian draws from NumPy seed 123.
At a state \((m,C)\), set \(A_C=\E\nabla^2V_\rho(X)\) and
\(B=C^{1/2}A_CC^{1/2}\).  The covariance update is
\[
  C^+=C^{1/2}\exp\!\left[
     \frac{h}{2\omega}
     \left\{-B+
       \frac{\omega+\tau\Tr B}{\omega+d\tau}\Id
     \right\}
  \right]C^{1/2},
  \qquad m^+=m-hC\E\nabla V_\rho(X).
\]
The dynamics expectations use one fixed scrambled Sobol design, transformed
coordinatewise by the Gaussian inverse CDF, with \(K=4096\) and seed 2026.
The same points are reused at every state and for every \((\omega,\tau)\) and
initialization.  The reference optimum uses an independent scrambled Sobol
design with \(K_{\rm ref}=8192\) and seed 2027, followed by Cholesky-coordinate
L-BFGS-B with a 2000-iteration budget, gradient tolerance \(10^{-6}\), and
function tolerance \(10^{-15}\).  The dynamics grid is
\[
  \omega\in\{0.25,0.5,1\},\quad
  \tau\in\{-\omega/(2d),0,\omega/(2d)\},\quad
  h=0.005,\quad T=40,
\]
with metrics saved every ten steps.  There are no random replicates: QMC is a
single common deterministic design once the seeds are fixed.

The five Gaussian initializations are
\begin{align*}
  \text{mean:}&\quad m_0=3\one/\sqrt d,\quad C_0=\Id,\\
  \text{volume-high/low:}&\quad m_0=0,\quad C_0=4\Id\ \text{or}\ \Id/4,\\
  \text{shape:}&\quad m_0=0,\quad
       C_0=\diag(e^2,e^{-2},1,\ldots,1),\\
  \text{mixed:}&\quad m_0=2\one/\sqrt d,\quad
       C_0=2\diag(e^{1.5},e^{-1.5},1,\ldots,1).
\end{align*}
For the non-Gaussian target the same formulas are applied in the
\((m_\star,C_\star)\)-whitened coordinates.  Thus, for example, the shape
covariance is
\(C_\star^{1/2}\diag(e^2,e^{-2},1,\ldots,1)C_\star^{1/2}\).

Both panels of \Cref{fig:num-affine-modes} show
\(T(\tau)/T(0)\) for every \(\omega\) and initialization.  The Gaussian panel
selects \(d=5\) and the first saved time at which the exact Gaussian KL,
normalized by its initial value, is at most \(10^{-4}\).  The non-Gaussian
panel uses the fixed-QMC objective gap, normalized by its initial gap, at
tolerance \(10^{-2}\).  A cell is left blank if either hitting time is not
finite within the horizon.

\paragraph{Data and figure lineage.}
\begingroup\small\raggedright
The configurations are \path{configs/omega_tau_modes/gaussian_target.yaml} and
\path{configs/omega_tau_modes/logconcave_target.yaml}; runners and kernels are
under \path{scripts/omega_tau_modes} and \path{src/omega_tau_modes}.  The report
builder reads the legacy committed locations
\path{outputs/gaussian_grid/summary.csv} and
\path{outputs/logconcave_grid/summary.csv}, through
\path{reports/make_report_assets.py::_fig_tau_speedup}.  The resulting assets
are \path{reports/assets/figs/fig_ot_gaussian_tau_speedup.pdf} and
\path{reports/assets/figs/fig_ot_logconcave_tau_speedup.pdf}.  To reproduce
these exact locations with the current runners, their \texttt{--outdir}
arguments must override the nested \texttt{output\_dir} values in the YAML
files.
\par\endgroup

\subsection{Non-log-concave boundary}

All selected runs are centered, scalar, deterministic, and use 160-point
Gauss--Hermite quadrature for
\[
  V_R(x)=\frac12x^2-2R^2\log\cosh(x/R),\qquad
  V_R''(x)=-1+2\tanh^2(x/R)\in[-1,1].
\]
For \(X\sim\N(0,c)\), write
\(A_R(c)=\E V_R''(X)\).  The plotted Fisher--Rao gradient squared is
\((1-cA_R(c))^2\).

The Riemannian cascade uses
\[
  c_{n+1}=c_n\exp\{h(1-c_nA_R(c_n))\},
  \quad R\in\{20,50,100,300,1000\},
  \quad c_0=1,\quad h=1.
\]
It has a 20-step cap and stops before the next update if
\(c_n>R^2/4\), if \(\log c\) would exceed 700, or on nonfinite arithmetic.
The first panel of \Cref{fig:num-nonconvex} plots every recorded iterate for all
five \(R\)'s.

The pole experiment uses the unprojected update
\[
  c_{n+1}=\frac{(1+h)c_n}{1+hc_nA_R(c_n)},
  \qquad R=1000,\quad h=1,\quad c_0=1-\varepsilon,
\]
with \(\varepsilon\in\{10^{-1},10^{-2},10^{-3},10^{-4},10^{-5},10^{-6}\}\)
and a two-step cap.  The middle panel plots the gradient squared after the first
step against \(\varepsilon\), together with the reference
\(4\varepsilon^{-2}\).

The clipped experiment uses
\[
  c_{n+1}=\operatorname{clip}_{[0.5,2]}
  \frac{(1+h)c_n}{1+hc_nA_R(c_n)},
  \qquad c_0=1,\qquad \beta=1.
\]
The report panel selects the primary safety fraction 0.5, hence
\[
  L_{\rm clip}=2\beta\lambda_+=4,\qquad
  h=0.5/L_{\rm clip}=0.125,
\]
and runs 300 steps for each of the five \(R\)'s.  The full output also contains
safety fractions \(0.25\) and \(0.9\), plus the deliberately outside-theorem
step \(h=1/(\beta\lambda_+)=0.5\); the selected panel excludes those rows.  At
step \(n\), the code
records
\begin{align*}
  D_n&=\KL(\N(0,c_n)\|\N(0,c_{n+1})),\\
  D_{\min}(N)&=\min_{0\leq n<N}D_n,\\
  B_N&=\frac{h}{N}\{\calE(0,c_0)-\calE(0,c_N)\}.
\end{align*}
The right panel plots \(D_{\min}(N)\) and \(B_N\).  In the recorded primary
runs, the upper clip first activates at step 2 for every \(R\); the step-2
displacement is nonzero, while the next and subsequent pinned steps have
\(D_n=0\).

\paragraph{Data and figure lineage.}
\begingroup\small\raggedright
The configuration is
\path{configs/natural_gradient_nonconvex_instability/nonconvex_instability.yaml};
the runner and kernels are in
\path{scripts/natural_gradient_nonconvex_instability} and
\path{src/natural_gradient_nonconvex_instability}.  The cascade panel reads
\path{outputs/natural_gradient_nonconvex_instability/results_long.csv} through
\path{plotting.py::fig_riemannian_cascade}; the pole panel reads
\path{outputs/natural_gradient_nonconvex_instability/kl_pole_summary.csv}
through \path{plotting.py::fig_kl_pole}; and the clipped panel reads
\path{outputs/natural_gradient_nonconvex_instability/clipped_kl_stationarity.csv}
through \path{plotting.py::fig_clipped_kl_stationarity}, restricted to
\texttt{dt\_rule=theorem\_safe}.  The assets are, respectively,
\path{reports/assets/figs/fig_nonconvex_riemannian_cascade.pdf},
\path{reports/assets/figs/fig_nonconvex_kl_pole.pdf}, and
\path{reports/assets/figs/fig_nonconvex_clipped_kl_stationarity.pdf}.
\par\endgroup

\subsection{Provenance limits of the committed snapshot}

\begingroup\small\raggedright
The scientific parameters above are recoverable from the configurations,
source, and raw columns used by the selected figures.  Some run-level
provenance files that the runners are designed to write are not committed:
\begin{enumerate}[label=(\roman*)]
  \item \path{outputs/natural_gradient_covariance_bootstrap} lacks
  \path{target_metadata.json} and \path{run_metadata.json};
  \item \path{outputs/natural_gradient_stl_variance} lacks
  \path{algorithm_results_long.csv}, \path{target_metadata.json},
  \path{run_metadata.json}, \path{estimator_target_metadata.json}, and
  \path{estimator_run_metadata.json}; and
  \item the Gaussian affine-mode runner does not write a run-metadata file.
\end{enumerate}
Consequently, the exact Python, NumPy, Torch, CUDA, and hardware versions of
those original runs cannot be reconstructed from this snapshot.  The
estimator-level STL CSV itself records the Torch/CUDA backend, but not the GPU
model; the algorithm-level backend is not recorded in the surviving summary
files.  The repository's \path{requirements.txt} and
\path{requirements-gpu.txt} give lower dependency bounds rather than a locked
environment.

These omissions do not change the selected data-to-figure mappings: the two
STL panels need only the surviving estimator and tail CSVs, and the covariance
panels need only the surviving long and transport-summary CSVs.  They do mean
that \path{reports/make_report_assets.py}, invoked with either
\texttt{--only stl\_variance} or
\texttt{--only covariance\_bootstrap}, does not run to completion from the
current snapshot, because the group-level builder also loads the missing
metadata or supplementary trajectory files before producing all of that
group's tables and figures.  Rerunning the corresponding experiment scripts
restores those files; no missing value has been guessed here.

General WFR experiments under \path{outputs/wfr_gradient_flow} are not used in
any figure or claim in this manuscript.
\par\endgroup
